\chapter{Assignment 6.1: Substituion Matrix}

I can start of making a comparation of what could be the substituion matrix for DNA and for proteins.

For example, in language, we can take an Italian text; If we decide to replace a word in italian with another in spanish (I lived in Italy, I can assure that is clearly true :)),
for example the test "Stasera andiamo in centro a cercare un po' di movida." movida is a synomyn of "nightlife" in english, "vita notturna" in italian.
The thing is that this swap does not break the meaning of the sentence, because the two words are similar enough to be understood in context.

Just like an Italian speaker understands "movida" because it shares a common Latin root and cultural context with "vita notturna," a protein can often continue to function if one hydrophobic amino acid is swapped for another similar hydrophobic one. 
Is here where a substitution matrix comes in handy, to turn this biological similarity into a mathematical score where a match gets a high positive value and a swap between similar amino acids gets a lower but still positive score, 
whereas swapping an amino acid for one with a completely different chemical personality, is like replacing a word in a sentence with a random gibberish sound that breaks the entire meaning. 
These matrices act as a lookup table that allows computer algorithms to recognize that two sequences are related even if they aren't identical, effectively measuring the "evolutionary cost" of every possible change. 
